\section{Desarrollo y Estado del Arte}
\label{sec:desarrollo}
Este es el núcleo del trabajo escrito y debe dividirse de forma lógica:

\subsection{Tecnologías de Fabricación Aditiva}
Descripción y comparación técnica entre la fusión selectiva por láser (SLM) y la fusión por haz de electrones (EBM). Discusión sobre la importancia de las atmósferas controladas (Argón/Alto vacío) para evitar la oxidación y fragilización (\textit{embrittlement}) del material.

\subsection{Evolución Microestructural y Postprocesado}
Análisis de las velocidades de enfriamiento extremas implícitas en los procesos de aditiva de lecho de polvo. Explicación detallada de la formación de la microestructura metaestable acicular (fase $\alpha'$) frente a la estructura bimodal convencional ($\alpha+\beta$). Justificación de la necesidad crítica de tratamientos térmicos de postprocesado para la relajación de tensiones residuales y la recuperación de la ductilidad del implante.

\subsection{Aplicaciones Clínicas e Interfase}
Soluciones desarrolladas por la comunidad científica utilizando estructuras porosas (andamios trabeculares) para optimizar el módulo elástico global y mejorar los ratios de osteointegración real en cirugía ortopédica y maxilofacial.
